\section{Data Preparation}
\label{sec:data-prep}

\subsection{Pipeline Architecture}

The data pipeline moves data from the source through staging, ingestion, and storage, ultimately reaching the data warehouse for analytics. The architecture is illustrated below:

\begin{verbatim}
Socrata API / CSV
    |
    v
PostgreSQL (Staging)
    |  * Schema validation
    |  * Bulk COPY FROM STDIN
    |  * Indexes on key columns
    |
    v
Sqoop (4 mappers, split-by id)
    |  * Parquet + Snappy (recommended)
    |
    v
HDFS (/user/team2/project/warehouse/)
    |
    v
Hive (team2_projectdb)
    |  * External table: crimes (Parquet on Sqoop output)
    |  * Managed table: crimes_optimized (partitioned + bucketed)
    |
    v
PySpark ML Pipeline (Stage III)
\end{verbatim}

\subsection{Data Ingestion}

\subsubsection{Data Collection}

The full dataset was downloaded from the Socrata SODA API using paginated requests of 500,000 rows per batch to prevent timeouts:

\begin{lstlisting}[language=bash, caption={Data collection script}]
bash scripts/data_collection.sh
\end{lstlisting}

Output: \texttt{data/chicago\_crimes\_raw.csv} ($\sim$1.9 GB) with MD5 checksum verification.

\subsubsection{PostgreSQL Staging}

The raw CSV was loaded into a PostgreSQL 14 database using \texttt{copy\_expert} for bulk ingestion. The \texttt{build\_db.py} script creates the \texttt{crimes} table with 22 columns, applies indexes on key columns (\texttt{arrest}, \texttt{year}, \texttt{primary\_type}, \texttt{date}, \texttt{district}, \texttt{community\_area}), and verifies row counts.

\subsubsection{Sqoop Import to HDFS}

Apache Sqoop was used to transfer data from PostgreSQL to HDFS with 4 mappers, split by the \texttt{id} column. Three format/codec combinations were benchmarked (Table~\ref{tab:sqoop-benchmark}).

\begin{table}[H]
\centering
\caption{Sqoop Format Benchmark}
\label{tab:sqoop-benchmark}
\begin{tabular}{lccc}
\toprule
\textbf{Format} & \textbf{Codec} & \textbf{HDFS Size} & \textbf{Import Time} \\
\midrule
Parquet & Snappy & \textbf{509 MB} & \textbf{54s} \\
AVRO & Snappy & 845 MB & 84s \\
Parquet & Gzip & 360 MB & 93s \\
\bottomrule
\end{tabular}
\end{table}

\textbf{Recommendation}: Parquet + Snappy was selected for its optimal balance of storage efficiency and import speed.

\subsection{Hive Optimization}

\subsubsection{Partitioned and Bucketed Table}

A managed table \texttt{crimes\_optimized} was created with the following optimisations:

\begin{itemize}
    \item \textbf{Partitioning}: By \texttt{year} (26 partitions: 2001--2026) — enables time-range pruning for queries.
    \item \textbf{Bucketing}: By \texttt{district} into 22 buckets — enables efficient join and group-by operations.
    \item \textbf{Storage}: Parquet + Snappy compression.
    \item \textbf{Files}: 370 total across all partitions.
\end{itemize}

The Hive DDL was executed as:

\begin{lstlisting}[language=SQL, caption={Hive table creation}]
CREATE TABLE team2_projectdb.crimes_optimized (
    id BIGINT, case_number STRING, date BIGINT,
    block STRING, iucr STRING, primary_type STRING,
    description STRING, location_description STRING,
    arrest BOOLEAN, domestic BOOLEAN, beat STRING,
    district INT, ward INT, community_area INT,
    fbi_code STRING, x_coordinate INT, y_coordinate INT,
    updated_on BIGINT, latitude DOUBLE, longitude DOUBLE,
    location STRING
)
PARTITIONED BY (year INT)
CLUSTERED BY (district) INTO 22 BUCKETS
STORED AS PARQUET
LOCATION '/user/team2/project/hive/warehouse/crimes_optimized';
\end{lstlisting}

\subsection{Data Cleaning}

\begin{table}[H]
\centering
\caption{Data Cleaning Operations}
\label{tab:cleaning}
\begin{tabularx}{\textwidth}{lX}
\toprule
\textbf{Issue} & \textbf{Handling} \\
\midrule
Date stored as epoch milliseconds & Converted to seconds; decomposed into hour, day\_of\_week, month via \texttt{from\_unixtime()} \\
Null categoricals & Filled with \texttt{"Unknown"} \\
Null district (47 rows) & Filled with -1 \\
Null community\_area (613,725 rows) & Filled with -1 \\
Missing coordinates (96,553 rows) & Excluded from ML training \\
Domestic flag & Encoded to numeric 1.0 / 0.0 \\
\bottomrule
\end{tabularx}
\end{table}

\subsection{Feature Engineering}

\subsubsection{Cyclical Time Encoding (SinCosTransformer)}

A custom PySpark ML transformer was implemented to encode cyclical temporal features into two continuous components, preserving circular relationships:

\begin{equation}
\text{sin\_component} = \sin\left(\frac{2\pi \times \text{value}}{\text{period}}\right),\quad
\text{cos\_component} = \cos\left(\frac{2\pi \times \text{value}}{\text{period}}\right)
\end{equation}

\begin{table}[H]
\centering
\caption{Cyclical Feature Encoding}
\label{tab:cyclical}
\begin{tabular}{lll}
\toprule
\textbf{Feature} & \textbf{Period} & \textbf{Output Columns} \\
\midrule
\texttt{hour\_of\_day} & 24 & \texttt{hour\_sin}, \texttt{hour\_cos} \\
\texttt{day\_of\_week} & 7 & \texttt{day\_sin}, \texttt{day\_cos} \\
\texttt{month} & 12 & \texttt{month\_sin}, \texttt{month\_cos} \\
\bottomrule
\end{tabular}
\end{table}

\subsubsection{Geospatial Encoding (GeoToECEFTransformer)}

A second custom transformer converts geodetic coordinates (latitude, longitude) to Earth-Centered Earth-Fixed (ECEF) coordinates using the WGS84 ellipsoid:

\begin{itemize}
    \item Semi-major axis ($a$): 6,378,137 m
    \item First eccentricity squared ($e^2$): $6.69437999014 \times 10^{-3}$
    \item Altitude: 0 (mean sea level)
\end{itemize}

Output: \texttt{ecef\_x}, \texttt{ecef\_y}, \texttt{ecef\_z} (in metres).

\subsubsection{Pipeline Stages}

The complete feature engineering pipeline consists of:

\begin{enumerate}
    \item SinCosTransformer $\times$ 3 (hour, day\_of\_week, month)
    \item GeoToECEFTransformer (lat, lon $\rightarrow$ ecef\_x, ecef\_y, ecef\_z)
    \item StringIndexer $\times$ 4 (primary\_type, iucr, location\_description, beat) with \texttt{handleInvalid="keep"}
    \item OneHotEncoder $\times$ 4
    \item VectorAssembler (all features $\rightarrow$ raw\_features vector)
    \item VarianceThresholdSelector (removes zero-variance features)
\end{enumerate}

The final feature vector includes:
\begin{itemize}
    \item OHE-encoded categoricals (4 groups)
    \item Numerical: district, community\_area, x\_coordinate, y\_coordinate, domestic
    \item Cyclical sin/cos: 6 columns
    \item ECEF geospatial: 3 columns
\end{itemize}
